\textbf{Phần 1. Trắc nghiệm lựa chọn}
\setcounterpageref{ex}{0}

\begin{ex} %[1D5V17]
Số lớn nhất có một chữ số là số nào?
\choice
{$8$}
{$10$}
{\True $9$}
{$7$}
\loigiai{
Số tự nhiên lớn nhất có một chữ số là số $9$.
\textbf{Đáp số:} $9$.\par
\textbf{Đáp án của câu hỏi là C}
}
\end{ex}

\begin{ex} %[1D5V17]
Số bé nhất có một chữ số là số nào?
\choice
{$1$}
{\True $0$}
{$2$}
{$3$}
\loigiai{
Số tự nhiên bé nhất có một chữ số là số $0$.
\textbf{Đáp số:} $0$.\par
\textbf{Đáp án của câu hỏi là B}
}
\end{ex}

\begin{ex} %[1D5V17]
Số bé nhất có hai chữ số là số nào?
\choice
{$0$}
{$20$}
{$9$}
{\True $10$}
\loigiai{
Số tự nhiên bé nhất có hai chữ số là số $10$.
\textbf{Đáp số:} $10$.\par
\textbf{Đáp án của câu hỏi là D}
}
\end{ex}

\begin{ex} %[1D5V17]
Số liền sau của số lớn nhất có một chữ số là số nào?
\choice
{\True $10$}
{$0$}
{$20$}
{$9$}
\loigiai{
Số lớn nhất có một chữ số là 9. Số liền sau của 9 là $9 + 1 = 10$.
\textbf{Đáp số:} $10$.\par
\textbf{Đáp án của câu hỏi là A}
}
\end{ex}

\begin{ex} %[1D5V17]
Số liền trước của số bé nhất có hai chữ số là số nào?
\choice
{$8$}
{$10$}
{\True $9$}
{$7$}
\loigiai{
Số bé nhất có hai chữ số là 10. Số liền trước của 10 là $10 - 1 = 9$.
\textbf{Đáp số:} $9$.\par
\textbf{Đáp án của câu hỏi là C}
}
\end{ex}

\begin{ex} %[1D5V17]
Có tất cả bao nhiêu số tự nhiên có một chữ số?
\choice
{$0$}
{\True $10$}
{$20$}
{$9$}
\loigiai{
Các số tự nhiên có một chữ số là: $0, 1, 2, 3, 4, 5, 6, 7, 8, 9$. Có tất cả $10$ số.
\textbf{Đáp số:} $10$.\par
\textbf{Đáp án của câu hỏi là B}
}
\end{ex}

\begin{ex} %[1D5V17]
Có bao nhiêu số tự nhiên có một chữ số khác 0?
\choice
{$8$}
{$10$}
{$7$}
{\True $9$}
\loigiai{
Các số đó là từ 1 đến 9, có tất cả $9$ số.
\textbf{Đáp số:} $9$.\par
\textbf{Đáp án của câu hỏi là D}
}
\end{ex}

\begin{ex} %[1D5V17]
Tìm số tự nhiên lớn hơn 7 và bé hơn 9.
\choice
{\True $8$}
{$7$}
{$9$}
{$6$}
\loigiai{
Số tự nhiên nằm giữa 7 và 9 là số $8$.
\textbf{Đáp số:} $8$.\par
\textbf{Đáp án của câu hỏi là A}
}
\end{ex}

\begin{ex} %[1D5V17]
Cho các số: $3, 8, 5, 10, 1$. Số lớn nhất trong các số đã cho là số nào?
\choice
{$0$}
{$20$}
{\True $10$}
{$9$}
\loigiai{
Trong các số đã cho, số lớn nhất là số $10$.
\textbf{Đáp số:} $10$.\par
\textbf{Đáp án của câu hỏi là C}
}
\end{ex}

\begin{ex} %[1D5V17]
Cho các số: $7, 2, 9, 4, 6$. Số bé nhất trong các số đã cho là số nào?
\choice
{$1$}
{\True $2$}
{$3$}
{$4$}
\loigiai{
Trong các số đã cho, số bé nhất là số $2$.
\textbf{Đáp số:} $2$.\par
\textbf{Đáp án của câu hỏi là B}
}
\end{ex}

\begin{ex} %[1D5V17]
Sắp xếp các số $4, 9, 2, 7, 5$ theo thứ tự từ lớn đến bé, số đứng ở vị trí thứ ba là số nào?
\choice
{$4$}
{$6$}
{$3$}
{\True $5$}
\loigiai{
Thứ tự từ lớn đến bé: $9 > 7 > 5 > 4 > 2$. Số đứng thứ ba là số $5$.
\textbf{Đáp số:} $5$.\par
\textbf{Đáp án của câu hỏi là D}
}
\end{ex}

\begin{ex} %[1D5V17]
Sắp xếp các số $10, 3, 6, 8, 0$ theo thứ tự từ bé đến lớn, số đứng ở vị trí thứ hai là số nào?
\choice
{\True $3$}
{$2$}
{$4$}
{$1$}
\loigiai{
Thứ tự từ bé đến lớn: $0 < 3 < 6 < 8 < 10$. Số đứng thứ hai là số $3$.
\textbf{Đáp số:} $3$.\par
\textbf{Đáp án của câu hỏi là A}
}
\end{ex}

\textbf{Phần 2. Trắc nghiệm đúng sai}
\setcounterpageref{ex}{0}

\begin{ex} %[1D5V17]
Xét tính đúng hoặc sai của các khẳng định toán học sau đây:
\choiceTF
{\True Kết quả của bài toán «Tìm số tự nhiên $x$ thỏa mãn điều kiện: $5 < x < 7$» là $6$.}
{Kết quả của bài toán «Có bao nhiêu số tự nhiên $y$ thỏa mãn điều kiện: $2 < y < 8$» là $4$.}
{Kết quả của bài toán «Hiệu của số lớn nhất có một chữ số và số bé nhất có một chữ số là bao nhiêu» là $8$.}
{\True Kết quả của bài toán «Tổng của số đứng liền sau số 3 và số đứng liền trước số 6 là bao nhiêu» là $9$.}
\loigiai{
\begin{itemize}
    \item Ý a) Đúng: Số tự nhiên lớn hơn 5 và bé hơn 7 là số $6$.
\textbf{Đáp số:} $6$.
    \item Ý b) Sai: Kết quả đúng là $5$, không phải $4$. Các số đó là: $3, 4, 5, 6, 7$. Có tất cả $5$ số.
\textbf{Đáp số:} $5$.
    \item Ý c) Sai: Kết quả đúng là $9$, không phải $8$. Hiệu là: $9 - 0 = 9$.
\textbf{Đáp số:} $9$.
    \item Ý d) Đúng: Số liền sau số 3 là 4. Số liền trước số 6 là 5.\nTổng của hai số là: $4 + 5 = 9$.
\textbf{Đáp số:} $9$.
\end{itemize}
\begin{itemize}
    \item a) \textbf{Đúng}.
    \item b) \textbf{Sai}.
    \item c) \textbf{Sai}.
    \item d) \textbf{Đúng}.
\end{itemize}
}
\end{ex}

\begin{ex} %[1D5V17]
Xét tính đúng hoặc sai của các bài toán và phép tính sau đây:
\choiceTF
{Kết quả của bài toán «Tìm hai số tự nhiên có tổng bằng 10 và hiệu bằng 6. Số lớn hơn là số mấy» là $9$.}
{\True Kết quả của bài toán «Hai số tự nhiên có tổng bằng 9 và hiệu bằng 3. Số bé hơn là số mấy» là $3$.}
{\True Trong dãy số tự nhiên từ 0 đến 10 có $6$ số chẵn.}
{Kết quả của bài toán «Hiệu của số lớn nhất có một chữ số và số bé nhất có một chữ số là bao nhiêu» là $8$.}
\loigiai{
\begin{itemize}
    \item Ý a) Sai: Kết quả đúng là $8$, không phải $9$. Hai số có tổng bằng 10 và hiệu bằng 6 là 8 và 2 (vì $8 + 2 = 10$ và $8 - 2 = 6$).\nSố lớn hơn là số $8$.
\textbf{Đáp số:} $8$.
    \item Ý b) Đúng: Hai số có tổng bằng 9 và hiệu bằng 3 là 6 và 3 (vì $6 + 3 = 9$ và $6 - 3 = 3$).\nSố bé hơn là số $3$.
\textbf{Đáp số:} $3$.
    \item Ý c) Đúng: Các số chẵn là: $0, 2, 4, 6, 8, 10$. Có tất cả $6$ số chẵn.
\textbf{Đáp số:} $6$.
    \item Ý d) Sai: Kết quả đúng là $9$, không phải $8$. Hiệu là: $9 - 0 = 9$.
\textbf{Đáp số:} $9$.
\end{itemize}
\begin{itemize}
    \item a) \textbf{Sai}.
    \item b) \textbf{Đúng}.
    \item c) \textbf{Đúng}.
    \item d) \textbf{Sai}.
\end{itemize}
}
\end{ex}

\textbf{Phần 3. Trả lời ngắn}
\setcounterpageref{ex}{0}

\begin{ex} %[1D5V17]
Trong dãy số tự nhiên từ 0 đến 10 có bao nhiêu số lẻ?
\shortans{5}
\loigiai{
\textbf{Đáp số:} $5$\par
Các số lẻ là: $1, 3, 5, 7, 9$. Có tất cả $5$ số lẻ.
\textbf{Đáp số:} $5$.
}
\end{ex}

\begin{ex} %[1D5V17]
Số 10 gồm mấy chục que tính?
\shortans{1}
\loigiai{
\textbf{Đáp số:} $1$\par
10 que tính gồm 1 chục que tính (ghi số 1).
\textbf{Đáp số:} $1$.
}
\end{ex}

\begin{ex} %[1D5V17]
Một chục quả trứng gà là bao nhiêu quả trứng gà?
\shortans{10}
\loigiai{
\textbf{Đáp số:} $10$\par
Một chục quả trứng tương ứng với $10$ quả trứng.
\textbf{Đáp số:} $10$.
}
\end{ex}

\begin{ex} %[1D5V17]
Lấy số lớn nhất có một chữ số trừ đi số đứng liền sau số 4. Kết quả bằng bao nhiêu?
\shortans{4}
\loigiai{
\textbf{Đáp số:} $4$\par
Số lớn nhất có một chữ số là 9. Số liền sau số 4 là 5: $9 - 5 = 4$.
\textbf{Đáp số:} $4$.
}
\end{ex}

\begin{ex} %[1D5V17]
Tìm số tự nhiên $a$ biết rằng: $a$ lớn hơn 8 và $a$ bé hơn hoặc bằng 10. Nếu $a$ là số lẻ thì $a$ bằng mấy?
\shortans{9}
\loigiai{
\textbf{Đáp số:} $9$\par
Các số tự nhiên lớn hơn 8 và không quá 10 là 9 và 10. Trong đó số lẻ là số $9$.
\textbf{Đáp số:} $9$.
}
\end{ex}

\begin{ex} %[1D5V17]
Có 10 que tính xếp thành các hình tam giác rời nhau (mỗi hình 3 que). Hỏi có thể xếp được nhiều nhất bao nhiêu hình tam giác?
\shortans{3}
\loigiai{
\textbf{Đáp số:} $3$\par
Mỗi hình tam giác cần 3 que: $3 + 3 + 3 = 9$ que tính. Vậy xếp được nhiều nhất $3$ hình tam giác (dư 1 que).
\textbf{Đáp số:} $3$.
}
\end{ex}
